\documentclass[11pt]{article}

\usepackage[margin=1.05in]{geometry}
\usepackage[T1]{fontenc}
\usepackage[utf8]{inputenc}
\usepackage{lmodern}
\usepackage{microtype}
\usepackage{amsmath,amssymb,amsthm,mathtools,mathrsfs,bm}
\usepackage{enumitem}
\usepackage{booktabs}
\usepackage{array}
\usepackage{hyperref}
\usepackage[nameinlink,capitalize]{cleveref}
\usepackage{xcolor}

\hypersetup{
  colorlinks=true,
  linkcolor=blue!45!black,
  citecolor=blue!45!black,
  urlcolor=blue!45!black,
  pdftitle={Archimedean support parity, shifted-Hurwitz defects, and compact CUE derivative moment stacks},
  pdfauthor={Project synthesis}
}

\allowdisplaybreaks

\newcommand{\C}{\mathbf C}
\newcommand{\R}{\mathbf R}
\newcommand{\Q}{\mathbf Q}
\newcommand{\Z}{\mathbf Z}
\newcommand{\N}{\mathbf N}
\newcommand{\B}{\mathbf B}
\newcommand{\F}{\mathbf F}
\newcommand{\PP}{\operatorname{PP}}
\newcommand{\Res}{\operatorname*{Res}}
\newcommand{\res}{\operatorname{res}}
\newcommand{\Tr}{\operatorname{Tr}}
\newcommand{\End}{\operatorname{End}}
\newcommand{\Hom}{\operatorname{Hom}}
\newcommand{\Spec}{\operatorname{Spec}}
\newcommand{\supp}{\operatorname{supp}}
\newcommand{\Haar}{\operatorname{Haar}}
\newcommand{\diag}{\operatorname{diag}}
\newcommand{\Span}{\operatorname{span}}
\newcommand{\id}{\operatorname{id}}
\newcommand{\rad}{\operatorname{rad}}
\newcommand{\ord}{\operatorname{ord}}
\newcommand{\coker}{\operatorname{coker}}
\newcommand{\im}{\operatorname{im}}
\newcommand{\kerop}{\operatorname{ker}}
\newcommand{\Lie}{\operatorname{Lie}}
\newcommand{\GL}{\operatorname{GL}}
\newcommand{\SL}{\operatorname{SL}}
\newcommand{\SU}{\operatorname{SU}}
\newcommand{\U}{\operatorname{U}}
\newcommand{\Mat}{\operatorname{Mat}}
\newcommand{\Part}{\mathcal P}
\newcommand{\eps}{\varepsilon}
\newcommand{\tauu}{\tau}
\newcommand{\ellp}{\ell}
\newcommand{\1}{\mathbf 1}
\newcommand{\dd}{\,d}
\newcommand{\wt}{\widetilde}
\newcommand{\wh}{\widehat}
\newcommand{\ot}{\otimes}
\newcommand{\op}{\mathrm{op}}
\newcommand{\fin}{\mathrm{fin}}
\newcommand{\func}{\mathrm{func}}
\newcommand{\diff}{\mathrm{diff}}
\newcommand{\vis}{\mathrm{vis}}
\newcommand{\hid}{\mathrm{hid}}
\newcommand{\HS}{\mathrm{HS}}
\newcommand{\SMod}{\mathrm{SMod}}
\newcommand{\Diag}{\mathrm{Diag}}
\newcommand{\canc}{\mathrm{canc}}
\newcommand{\sep}{\mathrm{sep}}
\newcommand{\idl}{\mathrm{idl}}
\newcommand{\pr}{\mathrm{pr}}
\newcommand{\sgn}{\operatorname{sgn}}

\newtheorem{theorem}{Theorem}[section]
\newtheorem{proposition}[theorem]{Proposition}
\newtheorem{lemma}[theorem]{Lemma}
\newtheorem{corollary}[theorem]{Corollary}
\newtheorem{definition}[theorem]{Definition}
\newtheorem{remark}[theorem]{Remark}
\newtheorem{construction}[theorem]{Construction}

\numberwithin{equation}{section}

\title{Archimedean support parity, shifted-Hurwitz defects, and compact CUE derivative moment stacks}
\author{Project synthesis}
\date{April 2026}

\begin{document}
\maketitle

\begin{abstract}
We construct a single algebraic-analytic framework joining four structures that have appeared separately in the split-zero globalization and compact CUE derivative-moment notes.  The first structure is the split-zero semiring
\[
  G(R)=R\sqcup\{\tau\},
\]
whose arithmetic reflection and Boolean support character separate the supported zero $0_R$ from the unsupported zero $\tau$.  We give a full theorem-level account of its support character, its semimodule classification by join-semilattice-indexed diagrams of $R$-modules, its cancellative reduction to $R$-modules, its Boolean-cubical free objects, and the conservative ideal and factorization consequences.  The second structure is the quadratic quotient/lift cover
\[
  u=s(s-1),\qquad v=2s-1,
\]
where function parity, differential parity, and residue parity differ.  We prove the residue reversal theorem and the projective inflow identity that places the finite two-end defect on the anti-invariant lift coefficient $v/u$ and cancels it at infinity.  The third structure is the shifted-Hurwitz transform
\[
  Z_{\mu,t}(s)=\int \zeta(s,1+ty)\,d\mu(y),
\]
whose completed principal part has quotient coefficient $1+t m_1(\mu)$ and lift coefficient $-t m_1(\mu)$.  We prove its all-orders zero-deformation recursion and its compact Bernstein-Hurwitz finite approximation stack.  The fourth structure is the CUE derivative-moment side: for
\[
  \Phi_N(z)=\det(zI-U),\qquad X_N=|\Phi_N'(1)|^2,
\]
we use the compact Mellin measure, the first moment formula, the $N=3$ phase-cover trace algebra, and the Grover--Mezzadri--Simm higher-derivative exponent.  We then prove the finite phase-space enlargement theorem: for every finite abelian group $A$ of order $r$, the nonidentity Weyl-Heisenberg sectors of $A\times\widehat A$ form an orthogonal basis of $\mathfrak{sl}_r(\C)$.  Consequently the first-derivative $s$-fold CUE exponent $s^2+2s=(s+1)^2-1$ is the dimension of the full traceless phase-space algebra for $r=s+1$, while the ordinary $A_{r-1}$ root lattice is only the diagonal augmentation skeleton.  All claims are proved from the displayed definitions and standard analytic facts about the Hurwitz and Dedekind zeta functions.
\end{abstract}

\tableofcontents

\section{Introduction and main results}

The paper of Grover--Mezzadri--Simm studies the joint CUE derivative moments
\begin{equation}\label{eq:GMS-moment}
  M_{\mu,\nu}(z,N)
  :=\mathbf E_{U(N)}\left[\prod_{i=1}^{\ell(\mu)}\Lambda_N^{(\mu_i)}(z)
  \prod_{j=1}^{\ell(\nu)}\overline{\Lambda_N^{(\nu_j)}(z)}\right],
  \qquad \Lambda_N(z)=\det(1-zU^\dagger).
\end{equation}
In the inner-disc regime their leading term is a sum over contingency tables; in the microscopic regime $z=1-c/N$ and $\ell(\mu)=\ell(\nu)=s$, the exponent of $N$ is
\begin{equation}\label{eq:GMS-exponent}
  |\mu|+|\nu|+s^2.
\end{equation}
For first derivative $s$-fold moments,
\begin{equation}
  \mu=\nu=(1^s),\qquad |\mu|=|\nu|=s,
\end{equation}
so \eqref{eq:GMS-exponent} becomes
\begin{equation}\label{eq:first-exp-sl}
  s^2+2s=(s+1)^2-1.
\end{equation}
Equation \eqref{eq:first-exp-sl} is not, by itself, a representation-theoretic theorem.  The point of this paper is to build the exact algebraic object whose dimension is the right-hand side: the full finite phase-space operator basis of $\mathfrak{sl}_{s+1}(\C)$, not merely the diagonal $A_s$ root lattice.  This is done in \cref{sec:phase-space}.

The CUE side is coupled to three other structures.  The first is the split-zero semiring
\begin{equation}\label{eq:GR-def-intro}
  G(R)=R\sqcup\{\tau\},
\end{equation}
where $\tau$ is the additive zero of the semiring, while the ordinary ring zero
\begin{equation}
  e:=0_R\in R\subset G(R)
\end{equation}
is a supported element.  The semiring has two canonical maps:
\begin{equation}
  p_R:G(R)\to R,
  \qquad
  \chi_R:G(R)\to\B,
\end{equation}
where $p_R$ records arithmetic value and $\chi_R$ records support.  The support character is the unique Boolean character; in particular
\begin{equation}
  \chi_R(\tau)=0,
  \qquad
  \chi_R(r)=1\quad(r\in R),
  \qquad
  \chi_R(0_R)=1.
\end{equation}
Thus the ordinary arithmetic zero lies on the supported sheet, while $\tau$ is the unique unsupported zero.

The second structure is the quotient/lift double
\begin{equation}\label{eq:uv-intro}
  u=s(s-1),\qquad v=2s-1,
  \qquad v^2=1+4u.
\end{equation}
Functions decompose into the invariant quotient sector $\C(u)$ and anti-invariant lift sector $v\C(u)$.  The principal parts with simple poles at $0$ and $1$ have basis
\begin{equation}\label{eq:qell-intro}
  q=\frac1u=\frac1{s-1}-\frac1s,
  \qquad
  \ellp=\frac vu=\frac1{s-1}+\frac1s.
\end{equation}
The subtlety is that residues are residues of differentials.  Since the involution $s\mapsto 1-s$ sends $ds$ to $-ds$, residues reverse the function parity.  We prove
\begin{equation}
  \res(q\,ds)=-e_0+e_1,
  \qquad
  \res(\ellp\,ds)=e_0+e_1.
\end{equation}
Therefore the visible two-end residue sector is not the function-even pole $1/u$; it is the differential-even, function-odd lift pole $v/u$.

The third structure is the shifted Hurwitz average.  For a compactly supported probability measure $\mu$ on $[0,\infty)$, write
\begin{equation}
  m_n(\mu)=\int_0^\infty y^n\,d\mu(y),
  \qquad
  Z_{\mu,t}(s)=\int_0^\infty \zeta(s,1+ty)\,d\mu(y),
\end{equation}
with completion
\begin{equation}
  \Lambda_{\mu,t}(s)=\pi^{-s/2}\Gamma(s/2)Z_{\mu,t}(s).
\end{equation}
Its principal part at $\{0,1\}$ is
\begin{equation}\label{eq:main-pp-intro}
  \PP_{\{0,1\}}\Lambda_{\mu,t}
  =\bigl(1+t m_1(\mu)\bigr)\frac1u
  -t m_1(\mu)\frac vu.
\end{equation}
The finite endpoint residue defect is
\begin{equation}
  -2t m_1(\mu),
\end{equation}
and the projective residue at infinity is
\begin{equation}
  +2t m_1(\mu).
\end{equation}
Thus the finite defect is canceled globally on $\mathbf P^1_s$.

For the CUE derivative variable
\begin{equation}
  X_N(U)=|\Phi_N'(1)|^2,
  \qquad \Phi_N(z)=\det(zI-U),
\end{equation}
the fixed-$N$ compact Mellin note proves that $X_N$ pushes Haar measure to a compact probability measure $\nu_N$ on the exact interval
\begin{equation}
  [0,L_N],\qquad L_N=N^2 4^{N-1}.
\end{equation}
The peeled measure $Y_N=X_N/N^2$ has first moment
\begin{equation}\label{eq:CUE-first-intro}
  m_1(\mu_N)=\frac{\mathbf E X_N}{N^2}
  =\frac{(N+1)(2N+1)}{6N}.
\end{equation}
Therefore \eqref{eq:main-pp-intro} gives an exact CUE shifted-Hurwitz defect with primitive lift coefficient
\begin{equation}
  -t\frac{(N+1)(2N+1)}{6N}.
\end{equation}
For $N=3$, this is $-14t/9$.

The paper is organized as follows.  \Cref{sec:split-zero} proves the split-zero support theorems and semimodule classification.  \Cref{sec:residue} proves the quotient/lift residue reversal and projective inflow.  \Cref{sec:hurwitz} proves the shifted-Hurwitz principal part and zero deformation.  \Cref{sec:cue-stack} proves the compact CUE-Hurwitz approximation stack.  \Cref{sec:phase-space} proves the finite phase-space enlargement theorem and the $\mathfrak{sl}_r$ basis theorem.  \Cref{sec:CUE-exponent} relates the Grover--Mezzadri--Simm exponent to this phase-space algebra.  \Cref{sec:N3} records the quadratic trace-cover mechanism behind the first nontrivial fixed-$N$ CUE layer.  \Cref{sec:synthesis} gives the consolidated dictionary.

\section{The split-zero semiring and its linear algebra}\label{sec:split-zero}

Let $R$ be a nonzero commutative ring with identity.  Define
\begin{equation}
  S=G(R):=R\sqcup\{\tau\}.
\end{equation}
For $r,s\in R$, set
\begin{equation}\label{eq:addition-GR}
  r\oplus s=r+s,
  \qquad
  r\oplus\tau=\tau\oplus r=r,
  \qquad
  \tau\oplus\tau=\tau,
\end{equation}
where $r+s$ is ordinary addition in $R$.  Multiplication is
\begin{equation}\label{eq:multiplication-GR}
  rs=rs\quad(r,s\in R),
  \qquad
  r\tau=\tau r=\tau,
  \qquad
  \tau^2=\tau.
\end{equation}
Then the additive zero of $S$ is $\tau$, not $0_R$.  We write
\begin{equation}
  e:=0_R\in R\subset S.
\end{equation}

\begin{proposition}[Boolean-coordinate normal form]\label{prop:DR}
There is a canonical semiring isomorphism
\begin{equation}
  G(R)\xrightarrow{\sim}
  D_R:=\{(r,b)\in R\times\B: b=0\Rightarrow r=0_R\}
       =\{(0_R,0)\}\cup (R\times\{1\}),
\end{equation}
given by
\begin{equation}
  \tau\longmapsto (0_R,0),
  \qquad
  r\longmapsto (r,1)\quad(r\in R).
\end{equation}
The two coordinate projections are the arithmetic reflection
\begin{equation}
  p_R(\tau)=0_R,
  \qquad p_R(r)=r,
\end{equation}
and the support character
\begin{equation}
  \chi_R(\tau)=0,
  \qquad \chi_R(r)=1\quad(r\in R).
\end{equation}
\end{proposition}

\begin{proof}
The displayed map is bijective.  If $r,s\in R$, then
\begin{equation}
  (r,1)+(s,1)=(r+s,1),
\end{equation}
which matches $r\oplus s=r+s$.  Also
\begin{equation}
  (r,1)+(0,0)=(r,1),
\end{equation}
which matches $r\oplus\tau=r$, and
\begin{equation}
  (0,0)+(0,0)=(0,0),
\end{equation}
which matches $\tau\oplus\tau=\tau$.  Multiplication is componentwise, with Boolean multiplication in the second coordinate; hence
\begin{equation}
  (r,1)(s,1)=(rs,1),\quad (r,1)(0,0)=(0,0),\quad (0,0)^2=(0,0),
\end{equation}
which is exactly \eqref{eq:multiplication-GR}.  Thus the map is a semiring isomorphism.
\end{proof}

\begin{theorem}[Unique support exception]\label{thm:unique-boolean-character}
There is exactly one semiring homomorphism
\begin{equation}
  G(R)\longrightarrow\B.
\end{equation}
It is $\chi_R$.  In particular
\begin{equation}
  \chi_R(\tau)=0,
  \qquad
  \chi_R(0_R)=1,
  \qquad
  \chi_R(r)=1\quad(r\in R).
\end{equation}
\end{theorem}

\begin{proof}
Let $f:G(R)\to\B$ be a semiring homomorphism.  Since $\tau$ is the additive zero of $G(R)$, one has
\begin{equation}
  f(\tau)=0.
\end{equation}
Since $1_R$ is the multiplicative identity, one has
\begin{equation}
  f(1_R)=1.
\end{equation}
Inside the supported copy $R\subset G(R)$,
\begin{equation}
  0_R=1_R+(-1_R).
\end{equation}
Applying $f$ gives
\begin{equation}
  f(0_R)=f(1_R)\oplus f(-1_R)=1\oplus f(-1_R)=1,
\end{equation}
because $1\oplus b=1$ in the Boolean semiring for both $b=0$ and $b=1$.

Suppose $r\in R$ and $f(r)=0$.  Since $r0_R=0_R$ in $R$, we get
\begin{equation}
  1=f(0_R)=f(r0_R)=f(r)f(0_R)=0\cdot1=0,
\end{equation}
a contradiction.  Hence $f(r)=1$ for every $r\in R$.  Therefore $f=\chi_R$.
\end{proof}

\begin{corollary}[Support pole]\label{cor:support-pole}
The unsupported element $\tau$, not $0_R$, is the support-pole point.  The ordinary zero $0_R$ lies in the supported sheet and maps to Boolean $1$.
\end{corollary}

\begin{proof}
This is exactly the statement
\begin{equation}
  \chi_R^{-1}(0)=\{\tau\},
  \qquad
  \chi_R^{-1}(1)=R.
\end{equation}
Since $0_R\in R$, it is supported.
\end{proof}

\subsection{Semimodule classification}

We now record the full replacement for ordinary $R$-linear algebra.  Let $M$ be a left $S$-semimodule.  Define the supported-zero projector
\begin{equation}\label{eq:epsilonM}
  \eps_M:M\to M,
  \qquad
  \eps_M(m):=em.
\end{equation}

\begin{lemma}\label{lem:epsilon}
The map $\eps_M$ is additive, idempotent, and satisfies
\begin{equation}\label{eq:absorbing-proj}
  m+\eps_M(m)=m
  \qquad(m\in M).
\end{equation}
\end{lemma}

\begin{proof}
Additivity follows from distributivity:
\begin{equation}
  \eps_M(m+n)=e(m+n)=em+en=\eps_M(m)+\eps_M(n).
\end{equation}
Idempotence follows from $e^2=e$:
\begin{equation}
  \eps_M^2(m)=e(em)=e^2m=em.
\end{equation}
Finally,
\begin{equation}
  m+\eps_M(m)=1\cdot m+e\cdot m=(1\oplus e)m.
\end{equation}
Since $1\oplus e=1_R+0_R=1_R$ in the supported copy $R$, this equals $m$.
\end{proof}

\begin{proposition}[Support semilattice]\label{prop:support-semilattice}
The image
\begin{equation}
  L_M:=eM=\im(\eps_M)
\end{equation}
is an idempotent commutative monoid under the ambient addition of $M$.  Hence
\begin{equation}
  l\vee l':=l+l'
\end{equation}
defines a join-semilattice with least element $0_M$.
\end{proposition}

\begin{proof}
If $l=em$ and $l'=en$, then
\begin{equation}
  l+l'=em+en=e(m+n)\in eM.
\end{equation}
Moreover
\begin{equation}
  l+l=em+em=(e\oplus e)m=em=l,
\end{equation}
because $e\oplus e=0_R+0_R=0_R=e$ in $R$.  Commutativity is inherited from the semimodule addition.  The additive zero $0_M$ belongs to $eM$ because $e0_M=0_M$, and it is the least element for the induced join.
\end{proof}

For $l\in L_M$, define the fiber
\begin{equation}
  M_l:=\{m\in M:\eps_M(m)=l\}.
\end{equation}

\begin{theorem}[Fibers are $R$-modules and transport maps]\label{thm:fibers-modules}
For each $l\in L_M$:
\begin{enumerate}[label=\textup{(\arabic*)}]
\item $M_l$ is canonically an $R$-module;
\item the zero element of this module is $l$;
\item if $l\le l'$ in $L_M$, then
\begin{equation}
  T_{l,l'}:M_l\to M_{l'},
  \qquad
  T_{l,l'}(m)=m+l'
\end{equation}
is $R$-linear;
\item $T_{l,l}=\id_{M_l}$ and $T_{l',l''}\circ T_{l,l'}=T_{l,l''}$ whenever $l\le l'\le l''$.
\end{enumerate}
\end{theorem}

\begin{proof}
Fix $l\in L_M$.  If $m,n\in M_l$, then
\begin{equation}
  \eps_M(m+n)=\eps_M(m)+\eps_M(n)=l+l=l,
\end{equation}
so $M_l$ is closed under addition.  The element $l$ lies in $M_l$ because $\eps_M(l)=l$.  It is an additive zero inside the fiber, since \cref{lem:epsilon} gives
\begin{equation}
  m+l=m+\eps_M(m)=m.
\end{equation}

For $r\in R$ and $m\in M_l$,
\begin{equation}
  \eps_M(rm)=e(rm)=(er)m.
\end{equation}
Since $er=0_R r=0_R=e$ in $R$, this equals $em=l$.  Thus $R$ acts fiberwise.  Addition and multiplication of scalars are inherited from the semimodule axioms, and additive inverses in $M_l$ are given by multiplication by $-1_R$:
\begin{equation}
  m+(-1_R)m=(1_R+(-1_R))m=0_Rm=em=l.
\end{equation}
Thus $M_l$ is an $R$-module with zero $l$.

If $l\le l'$, then $l+l'=l'$.  For $m\in M_l$,
\begin{equation}
  \eps_M(m+l')=l+l'=l',
\end{equation}
so $T_{l,l'}$ maps $M_l$ to $M_{l'}$.  It is $R$-linear because
\begin{equation}
  T_{l,l'}(rm)=rm+l'=rm+rl'=r(m+l'),
\end{equation}
using that every supported scalar acts trivially on $L_M$: if $r\in R$ and $l'=en$, then $rl'=ren=en=l'$ because $re=e$.  Additivity follows from
\begin{align}
  T_{l,l'}(m+n)&=m+n+l'\nonumber\\
  &=m+l'+n+l'\nonumber\\
  &=T_{l,l'}(m)+_{M_{l'}}T_{l,l'}(n),
\end{align}
where the last addition is the module addition in the fiber $M_{l'}$.  Finally
\begin{equation}
  T_{l,l}(m)=m+l=m
\end{equation}
by \eqref{eq:absorbing-proj}, and for $l\le l'\le l''$,
\begin{equation}
  T_{l',l''}(T_{l,l'}(m))=(m+l')+l''=m+l''=T_{l,l''}(m).
\end{equation}
\end{proof}

\begin{definition}[Support diagrams]\label{def:support-diagram}
A support diagram over $R$ is a triple
\begin{equation}
  A=\bigl(L,\{A_l\}_{l\in L},\{\rho_{l,l'}\}_{l\le l'}\bigr)
\end{equation}
where $L$ is a join-semilattice with least element $0$, each $A_l$ is an $R$-module, each $\rho_{l,l'}:A_l\to A_{l'}$ is $R$-linear, $\rho_{l,l}=\id$, and
\begin{equation}
  \rho_{l',l''}\circ\rho_{l,l'}=\rho_{l,l''}
  \qquad(l\le l'\le l'').
\end{equation}
Morphisms are the evident join-preserving maps on support semilattices together with compatible natural transformations of module fibers.  The category is denoted $\Diag_R$.
\end{definition}

\begin{construction}[From a support diagram to an $S$-semimodule]\label{constr:diagram-to-module}
Given $A=(L,A_l,\rho_{l,l'})\in\Diag_R$, define
\begin{equation}
  M(A):=\bigsqcup_{l\in L}A_l.
\end{equation}
For $x\in A_l$ and $y\in A_{l'}$, put
\begin{equation}\label{eq:diagram-addition}
  x\boxplus y:=\rho_{l,l\vee l'}(x)+_{A_{l\vee l'}}\rho_{l',l\vee l'}(y)
  \in A_{l\vee l'}.
\end{equation}
The additive zero is the zero element of $A_0$.  For $r\in R$, set $r\cdot x:=rx\in A_l$ and $\tau\cdot x:=0_{A_0}$.
\end{construction}

\begin{theorem}\label{thm:diagram-module}
\Cref{constr:diagram-to-module} makes $M(A)$ into a left $G(R)$-semimodule.
\end{theorem}

\begin{proof}
Commutativity of $\boxplus$ follows from commutativity of module addition.  For associativity, take
\begin{equation}
  x\in A_l,
  \qquad y\in A_m,
  \qquad z\in A_n,
  \qquad u=l\vee m\vee n.
\end{equation}
Then
\begin{align}
  (x\boxplus y)\boxplus z
  &=\rho_{l\vee m,u}\bigl(\rho_{l,l\vee m}(x)+\rho_{m,l\vee m}(y)\bigr)+\rho_{n,u}(z)\nonumber\\
  &=\rho_{l,u}(x)+\rho_{m,u}(y)+\rho_{n,u}(z),
\end{align}
by linearity and functoriality of the transition maps.  The same expression is obtained from $x\boxplus(y\boxplus z)$.

If $x\in A_l$, then
\begin{equation}
  x\boxplus 0_{A_0}=\rho_{l,l}(x)+\rho_{0,l}(0)=x,
\end{equation}
so $0_{A_0}$ is the additive zero.

For scalars $r,s\in R$, the identities
\begin{equation}
  (r+s)x=rx\boxplus sx,
  \qquad
  (rs)x=r(sx),
  \qquad
  1x=x
\end{equation}
are the corresponding module identities in the fixed fiber $A_l$.  If $x\in A_l$ and $y\in A_m$, then
\begin{align}
  r(x\boxplus y)&=r\rho_{l,l\vee m}(x)+r\rho_{m,l\vee m}(y)\nonumber\\
  &=\rho_{l,l\vee m}(rx)+\rho_{m,l\vee m}(ry)
  =(rx)\boxplus(ry).
\end{align}
Finally $\tau\cdot(x\boxplus y)=0_{A_0}=(\tau x)\boxplus(\tau y)$, and
\begin{equation}
  (r\oplus\tau)x=rx=(rx)\boxplus 0_{A_0}=(rx)\boxplus(\tau x).
\end{equation}
Thus all semimodule axioms hold.
\end{proof}

\begin{theorem}[Semimodule classification]\label{thm:semimodule-classification}
The assignments
\begin{equation}
  M\longmapsto \bigl(L_M,\{M_l\}_{l\in L_M},\{T_{l,l'}\}_{l\le l'}\bigr)
\end{equation}
and
\begin{equation}
  A\longmapsto M(A)
\end{equation}
extend to mutually quasi-inverse equivalences
\begin{equation}\label{eq:SMod-Diag-equivalence}
  G(R)\text{-}\SMod\simeq \Diag_R.
\end{equation}
\end{theorem}

\begin{proof}
The extraction of $L_M$, the fibers $M_l$, and the transport maps $T_{l,l'}$ is justified by \cref{prop:support-semilattice,thm:fibers-modules}.  Conversely, \cref{thm:diagram-module} constructs an $S$-semimodule from any support diagram.

Functoriality is direct.  If $f:M\to N$ is $S$-linear, then
\begin{equation}
  \eps_N(f(m))=ef(m)=f(em)=f(\eps_M(m)),
\end{equation}
so $f$ carries each fiber $M_l$ into the corresponding fiber of $N$, and compatibility with transport follows from additivity:
\begin{equation}
  f(T_{l,l'}m)=f(m+l')=f(m)+f(l')=T_{f(l),f(l')}f(m).
\end{equation}
A morphism of support diagrams gives an additive and scalar-compatible map on the disjoint union because naturality gives
\begin{align}
  \eta_{l\vee l'}(\rho_{l,l\vee l'}(x)+\rho_{l',l\vee l'}(y))
  &=\rho'_{\phi(l),\phi(l)\vee\phi(l')}(\eta_l(x))
  +\rho'_{\phi(l'),\phi(l)\vee\phi(l')}(\eta_{l'}(y)).
\end{align}
Starting from $M$, the reconstructed semimodule has the same underlying set, since the fibers $M_l$ partition $M$.  Its addition is the original addition because if $m\in M_l$ and $n\in M_{l'}$, then
\begin{equation}
  m+n=(m+l\vee l')+_{M_{l\vee l'}}(n+l\vee l'),
\end{equation}
which is the diagram formula \eqref{eq:diagram-addition}.  Scalar compatibility is immediate.  Starting from a support diagram, the extracted support semilattice, fibers, and transition maps are the original ones by construction.  Hence the two functors are quasi-inverse.
\end{proof}

\begin{corollary}[Cancellative sector]\label{cor:cancellative-sector}
Restriction along $p_R:G(R)\to R$ induces an equivalence
\begin{equation}
  R\text{-}\mathrm{Mod}\simeq G(R)\text{-}\SMod_{\canc},
\end{equation}
where the right-hand side is the full subcategory of cancellative $G(R)$-semimodules.
\end{corollary}

\begin{proof}
Let $M$ be cancellative.  From \eqref{eq:absorbing-proj},
\begin{equation}
  m+\eps_M(m)=m=m+0_M.
\end{equation}
Cancellation gives $\eps_M(m)=0_M$ for all $m$.  Hence $L_M$ has one element and \cref{thm:semimodule-classification} identifies $M$ with a single $R$-module.  Conversely, if $A$ is an $R$-module, define an $S$-action by $r\cdot a=ra$ for $r\in R$ and $\tau\cdot a=0$.  Then the underlying additive monoid is an abelian group, hence cancellative.
\end{proof}

\begin{theorem}[Free semimodules are Boolean cubes]\label{thm:free-boolean-cube}
Let $F_n:=G(R)^n$.  The support diagram of $F_n$ is the Boolean lattice
\begin{equation}
  \Part(\{1,\dots,n\})
\end{equation}
with fiber over $I\subseteq\{1,
\dots,n\}$ canonically equal to $R^I$, and transport along $I\subseteq J$ given by extension by zero in the coordinates of $J\setminus I$.
\end{theorem}

\begin{proof}
For $x=(x_1,\dots,x_n)\in G(R)^n$, the support is precisely
\begin{equation}
  \supp(x)=\{j:x_j\in R\}.
\end{equation}
Coordinatewise addition joins supports by union, so
\begin{equation}
  L_{F_n}\simeq\Part(\{1,\dots,n\}).
\end{equation}
The fiber over $I$ consists of vectors whose coordinates in $I$ are supported and whose coordinates outside $I$ are $\tau$.  Projection to the supported coordinates identifies this fiber with $R^I$.  If $I\subseteq J$, adding the support idempotent of $J$ replaces the new coordinates in $J\setminus I$ by $0_R$, i.e. it extends by zero.
\end{proof}

\subsection{Conservative ideal and factorization consequences}

For an ideal $J\triangleleft R$, write
\begin{equation}
  I_J:=J\cup\{\tau\}\subseteq G(R).
\end{equation}

\begin{theorem}[Ideal multiplication is unchanged]\label{thm:ideal-mult}
For all ideals $J,K\triangleleft R$,
\begin{equation}
  I_J I_K=I_{JK}.
\end{equation}
If $R$ is an integral domain, then the ideal class monoid of the nontrivial ideals of $G(R)$ is canonically isomorphic to the ordinary ideal class monoid of $R$.
\end{theorem}

\begin{proof}
Every product of an element of $I_J$ with an element of $I_K$ is either $\tau$ or an element of $JK$, so $I_JI_K\subseteq I_{JK}$.  Conversely, $\tau\in I_JI_K$ and every element of $JK$ is a finite sum of products $ab$ with $a\in J$ and $b\in K$, hence belongs to $I_JI_K$.

Assume now that $R$ is a domain.  For nonzero ideals $J_1,J_2\subset R$, ordinary ideal-class equivalence is
\begin{equation}
  J_1\sim_R J_2
  \Longleftrightarrow
  \exists a,b\in R\setminus\{0\}: aJ_1=bJ_2.
\end{equation}
For the corresponding $G(R)$-ideals,
\begin{equation}
  aI_J=aJ\cup\{\tau\},
\end{equation}
so
\begin{equation}
  aI_{J_1}=bI_{J_2}\Longleftrightarrow aJ_1=bJ_2.
\end{equation}
Thus the ideal class monoids coincide.
\end{proof}

\begin{proposition}[The separated size weight is nonmultiplicative]\label{prop:nonmult-weight}
Assume every nonzero quotient $R/J$ is finite and put $N(J)=|R/J|$.  Then
\begin{equation}
  w(J):=N(J)+1
\end{equation}
is not multiplicative on nonzero ideals.  More precisely,
\begin{equation}
  w(JK)=N(J)N(K)+1,
  \qquad
  w(J)w(K)=N(J)N(K)+N(J)+N(K)+1.
\end{equation}
Consequently
\begin{equation}
  \zeta^{\sep}_{G(R)}(s):=\sum_{0\ne J\triangleleft R}(N(J)+1)^{-s}
\end{equation}
does not admit an Euler product indexed by prime ideals via this weight.
\end{proposition}

\begin{proof}
The displayed identities are immediate from multiplicativity of $N$ on ideals and the definition $w(J)=N(J)+1$.  Their difference is $N(J)+N(K)>0$, so $w$ is not multiplicative.  An Euler product indexed by prime ideals would require multiplicativity of the local weight on coprime ideal products, so this separated sheet leaves the classical Euler category.
\end{proof}

\begin{proposition}[Units and irreducibles]\label{prop:units-irred}
Assume $R$ is an integral domain.  Then:
\begin{enumerate}[label=\textup{(\arabic*)}]
\item the units of $G(R)$ are exactly the units of $R$;
\item every nonzero supported irreducible of $G(R)$ is exactly an irreducible of $R$, and conversely;
\item the supported zero $e=0_R$ is reducible and divisibility-prime.
\end{enumerate}
\end{proposition}

\begin{proof}
If $u\in G(R)$ is a unit, then $u\ne\tau$ because $\tau x=\tau\ne1$.  Hence $u\in R$ and its inverse is supported, so $u$ is a unit of $R$.  The converse is immediate.

Let $x\in R\setminus\{0\}$.  If $x=ab$ in $G(R)$, neither factor can be $\tau$, since otherwise $ab=\tau$.  Hence $a,b\in R$, and factorization of $x$ in $G(R)$ is exactly factorization in $R$.

The element $e$ is reducible because $e=e^2$ and $e$ is not a unit.  Suppose $e\mid xy$.  Then $xy=ec$ for some $c\in G(R)$, so $xy\in\{e,\tau\}$.  If both $x$ and $y$ were supported nonzero elements, then $xy$ would be a supported nonzero element of the domain $R$, contradiction.  Thus one of $x,y$ lies in $\{e,\tau\}$, so $e$ divides one of them.
\end{proof}

\section{Quotient/lift parity and residue reversal}\label{sec:residue}

Set
\begin{equation}
  u=s(s-1),\qquad v=2s-1.
\end{equation}
Then
\begin{equation}
  u(1-s)=u(s),\qquad v(1-s)=-v(s),\qquad v^2=1+4u.
\end{equation}
Let
\begin{equation}
  K=\C(u),\qquad L=\C(s)=K(v),
\end{equation}
with deck involution
\begin{equation}
  \sigma(s)=1-s,
  \qquad
  \sigma(v)=-v.
\end{equation}
Let
\begin{equation}
  \mathcal P:=\left\{\frac as+\frac b{s-1}:a,b\in\C\right\}
\end{equation}
be the space of principal parts with at most simple poles at $0$ and $1$.  Define
\begin{equation}
  q:=\frac1u=\frac1{s-1}-\frac1s,
  \qquad
  \ellp:=\frac vu=\frac1{s-1}+\frac1s.
\end{equation}
Then
\begin{equation}
  \sigma(q)=q,
  \qquad
  \sigma(\ellp)=-\ellp.
\end{equation}
Thus
\begin{equation}
  \mathcal P=\mathcal P^+_{\func}\oplus\mathcal P^-_{\func},
  \qquad
  \mathcal P^+_{\func}=\C q,
  \qquad
  \mathcal P^-_{\func}=\C\ellp.
\end{equation}
Let
\begin{equation}
  E=\C e_0\oplus\C e_1,
\end{equation}
and let $\tau_E$ be the endpoint transposition
\begin{equation}
  \tau_E(e_0)=e_1,
  \qquad
  \tau_E(e_1)=e_0.
\end{equation}
Define
\begin{equation}
  \res:\mathcal P\,ds\to E,
  \qquad
  \res(P(s)\,ds)=\Res_{s=0}P(s)\,ds\,e_0+
  \Res_{s=1}P(s)\,ds\,e_1.
\end{equation}

\begin{theorem}[Residue parity reversal]\label{thm:residue-reversal}
One has
\begin{equation}
  \res(q\,ds)=-e_0+e_1,
  \qquad
  \res(\ellp\,ds)=e_0+e_1.
\end{equation}
Consequently the residue map reverses the $\Z/2$-parity: the function-even pole $q$ gives a residue-odd vector, while the function-odd pole $\ellp$ gives a residue-even vector.
\end{theorem}

\begin{proof}
From
\begin{equation}
  q=\frac1{s-1}-\frac1s
\end{equation}
we obtain
\begin{equation}
  \Res_{s=0}(q\,ds)=-1,
  \qquad
  \Res_{s=1}(q\,ds)=1.
\end{equation}
Thus $\res(q\,ds)=-e_0+e_1$, which is anti-invariant under $\tau_E$.

From
\begin{equation}
  \ellp=\frac1{s-1}+\frac1s
\end{equation}
we obtain
\begin{equation}
  \Res_{s=0}(\ellp\,ds)=1,
  \qquad
  \Res_{s=1}(\ellp\,ds)=1.
\end{equation}
Thus $\res(\ellp\,ds)=e_0+e_1$, which is invariant under $\tau_E$.

The structural reason is that residues are residues of differentials, not of functions.  Under the deck involution,
\begin{equation}
  \sigma^*(ds)=d(1-s)=-ds.
\end{equation}
Therefore
\begin{equation}
  \sigma^*(q\,ds)=-q\,ds,
  \qquad
  \sigma^*(\ellp\,ds)=(-\ellp)(-ds)=\ellp\,ds.
\end{equation}
Thus differential parity is function parity multiplied by the orientation character.
\end{proof}

\begin{corollary}[Finite defect detects the function-odd lift]\label{cor:finite-defect-lift}
Define
\begin{equation}
  \mathfrak D_{\fin}(P):=\Res_{s=0}P(s)\,ds+\Res_{s=1}P(s)\,ds.
\end{equation}
Then
\begin{equation}
  \mathfrak D_{\fin}(q)=0,
  \qquad
  \mathfrak D_{\fin}(\ellp)=2.
\end{equation}
Hence the visible two-end residue sector $e_0+e_1$ corresponds to the function-anti-invariant lift pole $v/u$.
\end{corollary}

\begin{proof}
The assertion follows by summing the two endpoint residues in \cref{thm:residue-reversal}.
\end{proof}

\begin{theorem}[Projective inflow]\label{thm:projective-inflow}
The divisor of $u=s(s-1)$ on $\mathbf P^1_s$ is
\begin{equation}
  (u)=[0]+[1]-2[\infty].
\end{equation}
Moreover
\begin{equation}
  \ellp\,ds=\frac vu\,ds=d\log u.
\end{equation}
Therefore, for every $B\in\C$,
\begin{equation}
  \Res_{s=0}(B\ellp\,ds)+\Res_{s=1}(B\ellp\,ds)=2B,
\end{equation}
while
\begin{equation}
  \Res_{s=\infty}(B\ellp\,ds)=-2B.
\end{equation}
The total projective residue is zero.
\end{theorem}

\begin{proof}
Since $u=s(s-1)$ has simple zeros at $0$ and $1$ and a pole of order $2$ at infinity, its divisor is $[0]+[1]-2[\infty]$.  Also
\begin{equation}
  du=(2s-1)\,ds=v\,ds,
\end{equation}
so
\begin{equation}
  d\log u=\frac{du}{u}=\frac vu\,ds=\ellp\,ds.
\end{equation}
The residue of $d\log u$ at a point is the order of $u$ at that point.  Hence the residues are $1,1,-2$, and multiplication by $B$ gives the displayed identities.
\end{proof}

\section{Completed shifted-Hurwitz transforms and endpoint defects}\label{sec:hurwitz}

Let $\mu$ be a compactly supported probability measure on $[0,\infty)$ and let
\begin{equation}
  m_n(\mu)=\int_0^\infty y^n\,d\mu(y).
\end{equation}
For $t\ge0$, define
\begin{equation}
  Z_{\mu,t}(s)=\int_0^\infty\zeta(s,1+ty)\,d\mu(y)
\end{equation}
and
\begin{equation}
  \Lambda_{\mu,t}(s)=\pi^{-s/2}\Gamma(s/2)Z_{\mu,t}(s).
\end{equation}

\begin{theorem}[Residues of the completed shifted average]\label{thm:hurwitz-residues}
One has
\begin{equation}
  \Res_{s=1}\Lambda_{\mu,t}=1,
\end{equation}
and
\begin{equation}
  \Res_{s=0}\Lambda_{\mu,t}=-1-2t m_1(\mu).
\end{equation}
Consequently
\begin{equation}
  \Res_{s=0}\Lambda_{\mu,t}+\Res_{s=1}\Lambda_{\mu,t}=-2t m_1(\mu).
\end{equation}
\end{theorem}

\begin{proof}
For $a>0$,
\begin{equation}
  \zeta(s,a)=\frac1{s-1}+O(1)\quad(s\to1)
\end{equation}
with residue $1$.  Since $\mu$ is a probability measure, $Z_{\mu,t}$ has residue $1$ at $s=1$.  The gamma factor satisfies
\begin{equation}
  \pi^{-1/2}\Gamma(1/2)=1,
\end{equation}
so the first residue is $1$.

At $s=0$,
\begin{equation}
  \zeta(0,a)=\frac12-a.
\end{equation}
Thus
\begin{equation}
  Z_{\mu,t}(0)=\int\left(\frac12-(1+ty)\right)d\mu(y)=-\frac12-tm_1(\mu).
\end{equation}
Also
\begin{equation}
  \Gamma(s/2)=\frac2s+O(1),
  \qquad
  \pi^{-s/2}=1+O(s).
\end{equation}
Therefore
\begin{equation}
  \Res_{s=0}\Lambda_{\mu,t}=2Z_{\mu,t}(0)=-1-2tm_1(\mu).
\end{equation}
Adding the two residues gives the asserted finite defect.
\end{proof}

\begin{theorem}[Quotient/lift principal part]\label{thm:pp-decomp}
The principal part of $\Lambda_{\mu,t}$ at $\{0,1\}$ is
\begin{equation}\label{eq:pp-decomp}
  \PP_{\{0,1\}}\Lambda_{\mu,t}
  =\bigl(1+t m_1(\mu)\bigr)\frac1u
  -t m_1(\mu)\frac vu.
\end{equation}
Equivalently,
\begin{equation}
  \PP_{\{0,1\}}\Lambda_{\mu,t}=A_{\mu,t}q+B_{\mu,t}\ellp,
  \qquad
  A_{\mu,t}=1+t m_1(\mu),
  \qquad
  B_{\mu,t}=-t m_1(\mu).
\end{equation}
\end{theorem}

\begin{proof}
The principal part has the form
\begin{equation}
  P(s)=\frac as+\frac b{s-1}
\end{equation}
with
\begin{equation}
  a=-1-2t m_1(\mu),
  \qquad b=1
\end{equation}
by \cref{thm:hurwitz-residues}.  Since
\begin{equation}
  Aq+B\ellp
  =A\left(\frac1{s-1}-\frac1s\right)+B\left(\frac1{s-1}+\frac1s\right),
\end{equation}
we have
\begin{equation}
  a=-A+B,
  \qquad b=A+B.
\end{equation}
Solving gives
\begin{equation}
  A=\frac{b-a}{2}=1+t m_1(\mu),
  \qquad
  B=\frac{a+b}{2}=-t m_1(\mu).
\end{equation}
\end{proof}

\begin{corollary}[Lift coefficient equals finite defect half]\label{cor:lift-defect-half}
For the principal part \eqref{eq:pp-decomp}, the finite endpoint defect is
\begin{equation}
  2B_{\mu,t}=-2t m_1(\mu),
\end{equation}
and the projective infinity contribution is
\begin{equation}
  -2B_{\mu,t}=+2t m_1(\mu).
\end{equation}
\end{corollary}

\begin{proof}
This is \cref{cor:finite-defect-lift,thm:projective-inflow} with $B=B_{\mu,t}$.
\end{proof}

\begin{corollary}[Critical-line value]\label{cor:critical-line-value}
For $s=1/2+i\gamma$,
\begin{equation}
  u=-\left(\gamma^2+\frac14\right),
  \qquad
  v=2i\gamma,
\end{equation}
and hence
\begin{equation}
  \PP_{\{0,1\}}\Lambda_{\mu,t}\left(\frac12+i\gamma\right)
  =-\frac{1+t m_1(\mu)}{\gamma^2+1/4}
  +i\frac{2t m_1(\mu)\gamma}{\gamma^2+1/4}.
\end{equation}
\end{corollary}

\begin{proof}
Substitute $u=-(\gamma^2+1/4)$ and $v=2i\gamma$ into \eqref{eq:pp-decomp}.
\end{proof}

\subsection{Quadratic trace orthogonality}

Let $A=R\oplus R\eta$ be a quadratic algebra over a characteristic-zero commutative ring $R$ with
\begin{equation}
  \eta^2=\Delta\in R,
  \qquad
  \sigma(a+b\eta)=a-b\eta.
\end{equation}
Define
\begin{equation}
  \langle x,y\rangle_A=\frac12\Tr_{A/R}(x\sigma(y)).
\end{equation}

\begin{lemma}[Quadratic trace form]\label{lem:quadratic-trace}
For $x=a+b\eta$ and $y=c+d\eta$,
\begin{equation}
  \langle x,y\rangle_A=ac-bd\Delta.
\end{equation}
In particular $A^+=R$ and $A^-=R\eta$ are orthogonal.
\end{lemma}

\begin{proof}
We compute
\begin{equation}
  (a+b\eta)\sigma(c+d\eta)=(a+b\eta)(c-d\eta)
  =ac-bd\eta^2+(bc-ad)\eta.
\end{equation}
Since $\eta^2=\Delta$ and $\Tr_{A/R}(r+s\eta)=2r$, the result follows.
\end{proof}

Apply this to
\begin{equation}
  A_z=\R[u,v]/(v^2-1-4u),
  \qquad \sigma(v)=-v.
\end{equation}
On the critical line $s=1/2+i\gamma$,
\begin{equation}
  A_{z,\gamma}=\R[v]/(v^2+4\gamma^2).
\end{equation}
For $\gamma\ne0$, this is a real form of $\C$, with $v\mapsto2i\gamma$.

\begin{corollary}[Critical-line real/imaginary splitting]\label{cor:critical-orthogonality}
For $\gamma\ne0$,
\begin{equation}
  A_{z,\gamma}^+=\R,
  \qquad
  A_{z,\gamma}^-=\R v=i\R.
\end{equation}
Moreover
\begin{equation}
  \langle a+bv,a+bv\rangle=a^2+4\gamma^2b^2,
\end{equation}
and in particular
\begin{equation}
  \left\langle\frac1u,\frac vu\right\rangle=0.
\end{equation}
\end{corollary}

\begin{proof}
Here $\Delta=-4\gamma^2$, so \cref{lem:quadratic-trace} gives
\begin{equation}
  \langle a+bv,a+bv\rangle=a^2-b^2(-4\gamma^2)=a^2+4\gamma^2b^2.
\end{equation}
Orthogonality of the even and odd eigenspaces follows from the same lemma.
\end{proof}

\subsection{All-orders shifted-zeta deformation}

Let
\begin{equation}
  L(s)=\sum_{n\ge1}a_n n^{-s}
\end{equation}
be a Dirichlet series with abscissa of absolute convergence $\sigma_a$.  For $|t|<1$, define
\begin{equation}
  L_t(s):=\sum_{n\ge1}a_n(n+t)^{-s}.
\end{equation}
Let
\begin{equation}
  (Tf)(s)=f(s+1),
  \qquad
  D=sT.
\end{equation}

\begin{theorem}[Universal shifted flow]\label{thm:shifted-flow}
For $|t|<1$ and $\Re(s)>\sigma_a$,
\begin{equation}\label{eq:Lt-expansion}
  L_t(s)=\sum_{m\ge0}(-1)^m\frac{(s)_m}{m!}t^m L(s+m).
\end{equation}
Equivalently,
\begin{equation}
  L_t=e^{-tD}L.
\end{equation}
Moreover, for every integer $n\ge0$,
\begin{equation}\label{eq:Lt-derivative}
  \partial_t^n L_t(s)=(-1)^n(s)_n L_t(s+n).
\end{equation}
\end{theorem}

\begin{proof}
For $|u|<1$,
\begin{equation}
  (1+u)^{-s}=\sum_{m\ge0}(-1)^m\frac{(s)_m}{m!}u^m.
\end{equation}
Apply this with $u=t/n$.  Since $|t|<1$, one has $|t/n|<1$ for every $n\ge1$.  Thus
\begin{equation}
  (n+t)^{-s}=n^{-s}\left(1+\frac tn\right)^{-s}
  =\sum_{m\ge0}(-1)^m\frac{(s)_m}{m!}t^m n^{-s-m}.
\end{equation}
Multiplying by $a_n$ and summing over $n$ gives \eqref{eq:Lt-expansion}.  Since
\begin{equation}
  D^mL(s)=(s)_mL(s+m),
\end{equation}
by induction on $m$, \eqref{eq:Lt-expansion} is precisely $e^{-tD}L$.  Finally,
\begin{equation}
  \partial_t^n(n+t)^{-s}=(-1)^n(s)_n(n+t)^{-s-n},
\end{equation}
and summing gives \eqref{eq:Lt-derivative}.
\end{proof}

\begin{corollary}[Logarithmic jets]\label{cor:log-jets}
Assume $L_t(s)\ne0$.  Put
\begin{equation}
  A_n(t;s):=\frac{\partial_t^nL_t(s)}{L_t(s)}
  =(-1)^n(s)_n\frac{L_t(s+n)}{L_t(s)}.
\end{equation}
Then
\begin{equation}\label{eq:cumulant-partition}
  \partial_t^n\log L_t(s)
  =\sum_{\pi\in\Pi_n}(-1)^{|\pi|-1}(|\pi|-1)!\prod_{B\in\pi}A_{|B|}(t;s),
\end{equation}
where $\Pi_n$ is the set of set partitions of $\{1,\dots,n\}$.
\end{corollary}

\begin{proof}
Formula \eqref{eq:cumulant-partition} is the standard moment-cumulant formula applied to the nonvanishing function $F(t)=L_t(s)$, with moments $A_n=F^{(n)}/F$.  It follows by differentiating $\log F$ repeatedly and grouping terms by set partitions.
\end{proof}

\begin{theorem}[Non-Eulerianity for nonintegral shifts]\label{thm:non-eulerianity}
Assume $a_1\ne0$.  If $0<t<1$, then $L_t$ is not an ordinary Dirichlet series in the variable $n^{-s}$.  Consequently $L_t$ admits no classical Euler product.
\end{theorem}

\begin{proof}
Suppose
\begin{equation}
  L_t(s)=\sum_{n\ge1}b_n n^{-s}
\end{equation}
for $\Re(s)$ sufficiently large.  Let $n_0$ be the least integer with $b_{n_0}\ne0$.  Then as $\sigma\to+\infty$,
\begin{equation}
  \sum b_n n^{-\sigma}=b_{n_0}n_0^{-\sigma}(1+o(1)).
\end{equation}
On the other hand,
\begin{equation}
  L_t(\sigma)=\sum_{n\ge1}a_n(n+t)^{-\sigma}=a_1(1+t)^{-\sigma}(1+o(1)),
\end{equation}
because the first term dominates.  Hence $n_0=1+t$, impossible since $1+t\in(1,2)$ is not an integer.
\end{proof}

\section{Zero deformation for compact Hurwitz averages}\label{sec:zero-deformation}

Let $\mu$ be compactly supported on $[0,L]$.  For $|t|L<1$,
\begin{equation}\label{eq:hurwitz-average-expansion}
  Z_{\mu,t}(s)=\sum_{n=0}^\infty
  (-t)^n\frac{(s)_n}{n!}m_n(\mu)\zeta(s+n),
\end{equation}
initially away from the pole locus and then by meromorphic continuation.

\begin{theorem}[All-order motion of a simple zeta zero]\label{thm:zero-motion}
Let $\rho$ be a simple nontrivial zero of $\zeta(s)$.  Define
\begin{equation}
  F_n(s):=(-1)^n\frac{(s)_n}{n!}m_n(\mu)\zeta(s+n).
\end{equation}
Then each $F_n$ is holomorphic at $s=\rho$.  There exists $\epsilon>0$ and a unique holomorphic branch
\begin{equation}
  \rho_\mu(t)=\rho+\sum_{N\ge1}c_Nt^N,
  \qquad |t|<\epsilon,
\end{equation}
such that
\begin{equation}
  Z_{\mu,t}(\rho_\mu(t))=0.
\end{equation}
Let
\begin{equation}
  a_{n,j}:=\frac1{j!}F_n^{(j)}(\rho).
\end{equation}
Then $a_{0,1}=\zeta'(\rho)\ne0$, and
\begin{equation}\label{eq:zero-recursion}
  c_N=-\frac1{\zeta'(\rho)}
  \sum_{\substack{n,j\ge0\\(n,j)\ne(0,1)}}
  a_{n,j}
  [t^N]\left(t^n\left(\sum_{m=1}^{N-1}c_mt^m\right)^j\right).
\end{equation}
Only finitely many terms contribute.  In particular,
\begin{equation}\label{eq:c1-zero}
  c_1=m_1(\mu)\frac{\rho\zeta(\rho+1)}{\zeta'(\rho)}.
\end{equation}
Moreover
\begin{equation}\label{eq:c2-zero}
  c_2=-\frac{\frac12\zeta''(\rho)c_1^2+F_1'(\rho)c_1+F_2(\rho)}{\zeta'(\rho)},
\end{equation}
where
\begin{equation}
  F_1(s)=-s m_1(\mu)\zeta(s+1),
  \qquad
  F_2(s)=\frac{s(s+1)}2m_2(\mu)\zeta(s+2).
\end{equation}
\end{theorem}

\begin{proof}
For $|t|L<1$, expansion \eqref{eq:hurwitz-average-expansion} gives a holomorphic expansion near $(s,t)=(\rho,0)$ because $\rho+n\ne1$ for all $n\ge0$: if $\rho+n=1$, then $\rho=1-n$, not a nontrivial zero.  At $t=0$,
\begin{equation}
  Z_{\mu,0}(s)=\zeta(s),
  \qquad
  \partial_s Z_{\mu,0}(\rho)=\zeta'(\rho)\ne0.
\end{equation}
The analytic implicit function theorem gives a unique holomorphic branch $\rho_\mu(t)$.

Let
\begin{equation}
  w(t)=\rho_\mu(t)-\rho=\sum_{N\ge1}c_Nt^N.
\end{equation}
Then
\begin{equation}
  0=Z_{\mu,t}(\rho+w(t))=
  \sum_{n,j\ge0}a_{n,j}t^n w(t)^j.
\end{equation}
Since $a_{0,0}=\zeta(\rho)=0$ and $a_{0,1}=\zeta'(\rho)$, the coefficient of $t^N$ is
\begin{equation}
  \zeta'(\rho)c_N+
  \sum_{\substack{n,j\ge0\\(n,j)\ne(0,1)}}
  a_{n,j}[t^N](t^nw(t)^j).
\end{equation}
In the second term no $c_N$ occurs: if $j\ge2$, the coefficient of $w(t)^j$ of degree $N$ uses only $c_1,\dots,c_{N-1}$, and if $n\ge1$, the coefficient has degree $N-n<N$.  Solving gives \eqref{eq:zero-recursion}.  For $N=1$, only $F_1(\rho)$ contributes besides $\zeta'(\rho)c_1$, and
\begin{equation}
  F_1(\rho)=-\rho m_1(\mu)\zeta(\rho+1),
\end{equation}
which gives \eqref{eq:c1-zero}.  For $N=2$, the coefficient equation is
\begin{equation}
  0=\zeta'(\rho)c_2+\frac12\zeta''(\rho)c_1^2+F_1'(\rho)c_1+F_2(\rho),
\end{equation}
which gives \eqref{eq:c2-zero}.
\end{proof}

\section{Compact CUE Mellin measures and finite Hurwitz zero stacks}\label{sec:cue-stack}

Let $U\in\U(N)$ be Haar distributed and define
\begin{equation}
  \Phi_N(z)=\det(zI-U),
  \qquad
  X_N(U)=|\Phi_N'(1)|^2.
\end{equation}
Let
\begin{equation}
  \nu_N:=(X_N)_*\Haar_{\U(N)}.
\end{equation}

\begin{theorem}[Exact compact support]\label{thm:compact-support}
For every $N\ge2$, $\nu_N$ is a compactly supported probability measure and
\begin{equation}
  \supp(\nu_N)=[0,L_N],
  \qquad
  L_N=N^2 4^{N-1}.
\end{equation}
Consequently, for every $k\ge0$,
\begin{equation}
  M_N(k):=\mathbf E_{\U(N)}|\Phi_N'(1)|^{2k}
  =\int_0^{L_N}t^k\,d\nu_N(t).
\end{equation}
\end{theorem}

\begin{proof}
If $z_1,\dots,z_N$ are the eigenvalues of $U$, then
\begin{equation}
  \Phi_N'(1)=\sum_{j=1}^N\prod_{i\ne j}(1-z_i).
\end{equation}
Since $|1-z_i|\le2$ on the unit circle,
\begin{equation}
  |\Phi_N'(1)|\le N2^{N-1}.
\end{equation}
Thus $0\le X_N\le N^24^{N-1}=L_N$.  The upper endpoint is achieved at $U=-I_N$, for which $\Phi_N(z)=(z+1)^N$ and $|\Phi_N'(1)|=N2^{N-1}$.  The lower endpoint is achieved at $U=I_N$, for which $\Phi_N(z)=(z-1)^N$ and $\Phi_N'(1)=0$ when $N\ge2$.  Since $\U(N)$ is connected and $X_N$ is continuous, the image of $X_N$ is a connected compact subset of $\R$ containing both endpoints, hence the full interval.  Haar measure has full support, so the support of the pushforward is that interval.  The Mellin identity is the definition of pushforward measure.
\end{proof}

Normalize
\begin{equation}
  \wt\nu_N:=(x\mapsto x/L_N)_*\nu_N
\end{equation}
on $[0,1]$ and write
\begin{equation}
  c_{N,n}:=\int_0^1x^n\,d\wt\nu_N(x)=\frac{M_N(n)}{L_N^n}.
\end{equation}
For $m\ge0$ and $0\le j\le m$, define
\begin{equation}\label{eq:bernstein-weights}
  \lambda_{N,m,j}:=\binom mj(-1)^{m-j}\Delta^{m-j}c_{N,j}
  =\binom mj\int_0^1x^j(1-x)^{m-j}\,d\wt\nu_N(x),
\end{equation}
where $\Delta c_j=c_{j+1}-c_j$.

\begin{theorem}[Bernstein-Hausdorff reconstruction]\label{thm:bernstein}
The weights \eqref{eq:bernstein-weights} satisfy
\begin{equation}
  \lambda_{N,m,j}\ge0,
  \qquad
  \sum_{j=0}^m\lambda_{N,m,j}=1.
\end{equation}
The atomic measures
\begin{equation}
  \wt\nu_{N,m}:=\sum_{j=0}^m\lambda_{N,m,j}\delta_{j/m}
\end{equation}
converge weakly to $\wt\nu_N$.  Hence the integer moment lattice $\{M_N(n)\}_{n\ge0}$ constructively determines the whole compact measure and therefore $M_N(k)$ for every real $k\ge0$.
\end{theorem}

\begin{proof}
The equality in \eqref{eq:bernstein-weights} follows from
\begin{equation}
  (-1)^r\Delta^r c_{N,j}=\int_0^1x^j(1-x)^r\,d\wt\nu_N(x),
\end{equation}
proved by induction on $r$.  Positivity follows because the integrand is nonnegative.  Summing over $j$ gives
\begin{equation}
  \sum_{j=0}^m\lambda_{N,m,j}
  =\int_0^1\sum_{j=0}^m\binom mjx^j(1-x)^{m-j}\,d\wt\nu_N(x)=1.
\end{equation}
For every continuous $f$ on $[0,1]$,
\begin{equation}
  \int f\,d\wt\nu_{N,m}
  =\int_0^1\sum_{j=0}^m f(j/m)\binom mjx^j(1-x)^{m-j}\,d\wt\nu_N(x).
\end{equation}
The inner sum is the Bernstein polynomial $B_mf(x)$, which converges uniformly to $f(x)$ on $[0,1]$.  Therefore the integrals converge.  Taking $f(x)=x^k$ gives the positive real moments after rescaling by $L_N^k$.
\end{proof}

Define the finite Hurwitz approximant
\begin{equation}\label{eq:finite-Hurwitz-approx}
  Z^{[m]}_{N,t}(s):=\sum_{j=0}^m\lambda_{N,m,j}\,
  \zeta\left(s,1+tL_N\frac jm\right).
\end{equation}

\begin{theorem}[Finite Bernstein-Hurwitz zero stack]\label{thm:finite-zero-stack}
For fixed $N$ and $t\ge0$,
\begin{equation}
  Z^{[m]}_{N,t}(s)\longrightarrow
  Z_{N,t}(s):=\int_0^{L_N}\zeta(s,1+tx)\,d\nu_N(x)
\end{equation}
locally uniformly on compact subsets of $\C\setminus\{1\}$.  Consequently, if $\Omega\subset\C\setminus\{1\}$ is bounded with piecewise smooth boundary and $Z_{N,t}$ has no zeros on $\partial\Omega$, then for all sufficiently large $m$, the number of zeros of $Z^{[m]}_{N,t}$ in $\Omega$, counted with multiplicity, equals the number of zeros of $Z_{N,t}$ in $\Omega$.
\end{theorem}

\begin{proof}
Let
\begin{equation}
  f_s(x)=\zeta(s,1+tL_Nx),\qquad x\in[0,1].
\end{equation}
For every compact $K\subset\C\setminus\{1\}$, the function $(s,x)\mapsto f_s(x)$ is continuous in $x$, holomorphic in $s$, and uniformly bounded on $K\times[0,1]$.  Weak convergence $\wt\nu_{N,m}\Rightarrow\wt\nu_N$ gives uniform convergence of the integrals on $K$, since the family $\{f_s:s\in K\}$ is equicontinuous on $[0,1]$ by compactness.  The zero-counting assertion follows from Hurwitz's theorem or from Rouche's theorem on $\partial\Omega$.
\end{proof}

\begin{theorem}[Exact first CUE derivative moment]\label{thm:first-CUE-moment}
For every $N\ge1$,
\begin{equation}
  \mathbf E_{\U(N)}|\Phi_N'(1)|^2
  =\sum_{m=1}^N m^2
  =\frac{N(N+1)(2N+1)}6.
\end{equation}
Consequently the peeled variable $Y_N=X_N/N^2$ has first moment
\begin{equation}
  m_1(\mu_N)=\frac{(N+1)(2N+1)}{6N}.
\end{equation}
\end{theorem}

\begin{proof}
Write
\begin{equation}
  \Phi_N(z)=\sum_{j=0}^N(-1)^j e_j(U)z^{N-j},
\end{equation}
where $e_j(U)$ is the $j$-th elementary symmetric polynomial in the eigenvalues of $U$.  Then
\begin{equation}
  \Phi_N'(1)=\sum_{j=0}^{N-1}(-1)^j(N-j)e_j(U).
\end{equation}
The functions $e_j$ are irreducible exterior-power characters of $\U(N)$ and are orthonormal under Haar measure:
\begin{equation}
  \mathbf E[e_j(U)\overline{e_k(U)}]=\delta_{jk}.
\end{equation}
Therefore
\begin{equation}
  \mathbf E|\Phi_N'(1)|^2=\sum_{j=0}^{N-1}(N-j)^2=\sum_{m=1}^Nm^2.
\end{equation}
The closed formula for the sum of squares is standard, and division by $N^2$ gives the peeled first moment.
\end{proof}

\begin{corollary}[CUE shifted-Hurwitz principal part]\label{cor:CUE-pp}
For the peeled CUE measure $\mu_N=(Y_N)_*\Haar$,
\begin{equation}
  \PP_{\{0,1\}}\Lambda_{N,t}
  =\left(1+t\frac{(N+1)(2N+1)}{6N}\right)\frac1u
  -t\frac{(N+1)(2N+1)}{6N}\frac vu.
\end{equation}
For $N=3$ this is
\begin{equation}
  \PP_{\{0,1\}}\Lambda_{3,t}
  =\left(1+\frac{14t}{9}\right)\frac1u-\frac{14t}{9}\frac vu.
\end{equation}
\end{corollary}

\begin{proof}
Apply \cref{thm:pp-decomp} with the first moment from \cref{thm:first-CUE-moment}.  For $N=3$ the value is $(4)(7)/(18)=14/9$.
\end{proof}

\section{Finite phase space and the traceless operator algebra}\label{sec:phase-space}

Let $A$ be a finite abelian group of order $r$.  Let
\begin{equation}
  \widehat A=\Hom(A,\C^\times),
  \qquad
  H_A=\C[A].
\end{equation}
For $a\in A$ and $\chi\in\widehat A$, define
\begin{equation}
  T_a e_x=e_{a+x},
  \qquad
  M_\chi e_x=\chi(x)e_x,
  \qquad
  W_{a,\chi}=T_aM_\chi.
\end{equation}

\begin{lemma}[Weyl product]\label{lem:weyl-product}
One has
\begin{equation}
  W_{a,\chi}W_{b,\psi}=\chi(b)W_{a+b,\chi\psi}.
\end{equation}
Consequently
\begin{equation}
  [W_{a,\chi},W_{b,\psi}]
  =\bigl(\chi(b)-\psi(a)\bigr)W_{a+b,\chi\psi}.
\end{equation}
\end{lemma}

\begin{proof}
Since
\begin{equation}
  M_\chi T_b e_x=M_\chi e_{b+x}=\chi(b+x)e_{b+x}=\chi(b)T_bM_\chi e_x,
\end{equation}
we have $M_\chi T_b=\chi(b)T_bM_\chi$.  Therefore
\begin{equation}
  T_aM_\chi T_bM_\psi=\chi(b)T_{a+b}M_{\chi\psi}.
\end{equation}
The commutator formula follows by subtracting the same expression with $(a,\chi)$ and $(b,\psi)$ exchanged.
\end{proof}

\begin{theorem}[Finite phase-space basis theorem]\label{thm:phase-space-basis}
The operators
\begin{equation}
  \{W_{a,\chi}:a\in A,\ \chi\in\widehat A\}
\end{equation}
form an orthogonal basis of $\End(H_A)$ for the Hilbert-Schmidt form
\begin{equation}
  \langle S,T\rangle_{\HS}=\Tr(S^*T).
\end{equation}
More precisely,
\begin{equation}
  \Tr(W_{a,\chi}^*W_{b,\psi})=r\,\delta_{a,b}\delta_{\chi,\psi}.
\end{equation}
Moreover $W_{0,1}=I$, and every other $W_{a,\chi}$ is traceless.  Hence
\begin{equation}\label{eq:phase-sl}
  \Span_\C\{W_{a,\chi}:(a,\chi)\ne(0,1)\}
  =\mathfrak{sl}(H_A)\cong\mathfrak{sl}_r(\C).
\end{equation}
\end{theorem}

\begin{proof}
For the trace, observe that
\begin{equation}
  W_{a,\chi}e_x=\chi(x)e_{a+x}.
\end{equation}
This contributes to the trace only if $a=0$.  If $a=0$, then
\begin{equation}
  \Tr(W_{0,\chi})=\sum_{x\in A}\chi(x),
\end{equation}
which equals $r$ if $\chi=1$ and $0$ otherwise by character orthogonality.

Since $W_{a,\chi}^*W_{b,\psi}$ is a scalar multiple of $W_{b-a,\psi\chi^{-1}}$, the same trace computation gives the orthogonality formula.  There are $r^2$ such operators and $\dim\End(H_A)=r^2$, so they form a basis.  The identity is $W_{0,1}$.  All other basis elements have trace zero, and the traceless endomorphisms have dimension $r^2-1$, proving \eqref{eq:phase-sl}.
\end{proof}

For $A=\Z/r\Z$, set
\begin{equation}
  \omega=e^{2\pi i/r},
  \qquad
  Xe_j=e_{j+1},
  \qquad
  Ze_j=\omega^j e_j.
\end{equation}
Then $ZX=\omega XZ$ and
\begin{equation}
  W_{a,b}=X^aZ^b,
  \qquad a,b\in\Z/r\Z.
\end{equation}
Let $E_{ij}$ denote the elementary matrix sending $e_j$ to $e_i$.  Then
\begin{equation}\label{eq:XaZb-elementary}
  X^aZ^b=\sum_{j\in\Z/r\Z}\omega^{bj}E_{j+a,j}.
\end{equation}

\begin{theorem}[Fourier root-space decomposition]\label{thm:fourier-root}
There is an orthogonal vector-space decomposition
\begin{equation}
  \mathfrak{sl}_r(\C)=\mathfrak h\oplus\bigoplus_{a\in\Z/r\Z\setminus\{0\}}V_a,
\end{equation}
where
\begin{equation}
  \mathfrak h=\Span_\C\{Z^b:1\le b\le r-1\}
\end{equation}
is the diagonal traceless Cartan subalgebra and
\begin{equation}
  V_a=\Span_\C\{X^aZ^b:b\in\Z/r\Z\}
  =\bigoplus_{j\in\Z/r\Z}\C E_{j+a,j}.
\end{equation}
Thus the nonzero displacement sectors are Fourier transforms of fixed-displacement root-space slices.
\end{theorem}

\begin{proof}
For $a=0$,
\begin{equation}
  Z^b=\sum_j\omega^{bj}E_{jj}.
\end{equation}
The nontrivial Fourier characters $b=1,\dots,r-1$ span the hyperplane of diagonal matrices whose trace is zero, hence span $\mathfrak h$.

For $a\ne0$, equation \eqref{eq:XaZb-elementary} shows
\begin{equation}
  X^aZ^b=\sum_j\omega^{bj}E_{j+a,j}.
\end{equation}
As $b$ varies, the coefficient matrix $(\omega^{bj})_{b,j}$ is the invertible Fourier matrix.  Hence the $X^aZ^b$ span exactly the elementary matrices $E_{j+a,j}$.  The dimensions add to
\begin{equation}
  (r-1)+r(r-1)=r^2-1.
\end{equation}
\end{proof}

\begin{corollary}[Where $A_{r-1}$ sits]\label{cor:A-root-lattice}
The integral diagonal augmentation lattice
\begin{equation}
  \Z^r_0=\{(n_0,
\dots,n_{r-1})\in\Z^r:\sum_jn_j=0\}
\end{equation}
is the $A_{r-1}$ root lattice.  Its complexification is the Cartan subalgebra $\mathfrak h$.  The full finite phase-space double $A\times\widehat A$ is not the $A_{r-1}$ root lattice; its nonidentity Weyl sectors form a basis of the entire traceless matrix algebra $\mathfrak{sl}_r(\C)$.
\end{corollary}

\begin{proof}
The root lattice of $\mathfrak{sl}_r$ is generated by $e_i-e_j$ inside $\Z^r_0$, and the simple roots $e_i-e_{i+1}$ give the Cartan matrix of type $A_{r-1}$.  Its complexification is the traceless diagonal subalgebra.  The second assertion is \cref{thm:phase-space-basis,thm:fourier-root}.
\end{proof}

\section{CUE derivative exponents and the finite phase-space enlargement}\label{sec:CUE-exponent}

Let $\mu,\nu$ be lists of nonnegative integers and let $M_{\mu,\nu}(z,N)$ be \eqref{eq:GMS-moment}.  In the microscopic/unit-circle regime of Grover--Mezzadri--Simm, when $\ell(\mu)=\ell(\nu)=s$, the exponent of $N$ is
\begin{equation}
  |\mu|+|\nu|+s^2.
\end{equation}
For first derivatives,
\begin{equation}
  \mu=\nu=(1^s),
  \qquad |\mu|=|\nu|=s.
\end{equation}
Thus the exponent is
\begin{equation}
  s^2+2s=(s+1)^2-1.
\end{equation}

\begin{theorem}[First-jet block realization]\label{thm:first-jet-block}
Let $W=\C^s$ and define
\begin{equation}
  \mathfrak B_s:=\End(W)\oplus W\oplus W^\vee.
\end{equation}
The map
\begin{equation}
  \Theta:\mathfrak B_s\to\mathfrak{sl}(W\oplus\C),
  \qquad
  \Theta(A,u,\varphi)=
  \begin{pmatrix}
    A & u\\
    \varphi & -\Tr(A)
  \end{pmatrix}
\end{equation}
is a vector-space isomorphism.  Consequently
\begin{equation}
  \dim\mathfrak B_s=s^2+2s=\dim\mathfrak{sl}_{s+1}(\C).
\end{equation}
\end{theorem}

\begin{proof}
Every element of $\mathfrak{sl}(W\oplus\C)$ has a unique block form
\begin{equation}
  \begin{pmatrix}
    B & u\\
    \varphi & \lambda
  \end{pmatrix},
  \qquad
  B\in\End(W),\ u\in W,
  \ \varphi\in W^\vee,
\end{equation}
with trace condition $\Tr(B)+\lambda=0$.  Hence $\lambda=-\Tr(B)$ and the displayed map is bijective.  The dimension is $s^2+s+s=s^2+2s=(s+1)^2-1$.
\end{proof}

\begin{corollary}[Phase-space model for the exponent]\label{cor:CUE-phase-exponent}
Let $A=\Z/(s+1)\Z$.  Then
\begin{equation}
  |A\times\widehat A|-1=(s+1)^2-1=s^2+2s.
\end{equation}
The nonidentity phase-space sectors form an orthogonal basis of $\mathfrak{sl}_{s+1}(\C)$.
\end{corollary}

\begin{proof}
This is \cref{thm:phase-space-basis} with $r=s+1$.
\end{proof}

\begin{remark}
\Cref{cor:CUE-phase-exponent} is a dimension-level theorem and an operator-basis theorem.  It does not assert that the Kostka determinant formula is literally a character formula for an $\SL_{s+1}$-module.  That stronger assertion would require a functorial identification of the determinant/Kostka expression with a representation-theoretic object.  What is proved here is the exact finite phase-space enlargement behind the exponent.
\end{remark}

\section{The first nontrivial CUE trace cover at $N=3$}\label{sec:N3}

The compact Mellin peeling note develops the first nontrivial finite-$N$ layer at $N=3$.  We record the algebraic trace-cover part because it is the exact CUE analogue of the quotient/lift splitting from \cref{sec:residue}.

Let
\begin{equation}
  u=|\lambda|^2,
  \qquad
  v=|c|^2,
  \qquad
  w=|d|^2=\frac{1-u+2v}{2},
\end{equation}
and let
\begin{equation}
  z=c\overline d\,\overline\lambda.
\end{equation}
The first peeled algebra gives
\begin{equation}\label{eq:Ruv}
  z+\overline z=R(u,v):=\frac{9u^2-16uv-10u-16v+1}{16},
\end{equation}
and
\begin{equation}\label{eq:Puv}
  z\overline z=P(u,v):=uvw=uv\frac{1-u+2v}{2}.
\end{equation}
Hence $z$ satisfies
\begin{equation}
  T^2-R(u,v)T+P(u,v)=0.
\end{equation}
Define
\begin{equation}
  J:=z-\frac R2.
\end{equation}
Then $\overline J=-J$ and
\begin{equation}
  J^2=-P+\frac{R^2}{4}=-S(u,v),
  \qquad
  S(u,v):=P(u,v)-\frac{R(u,v)^2}{4}.
\end{equation}
Put
\begin{equation}
  X:=9u-16v+1,
  \qquad
  \Delta:=36u-X^2.
\end{equation}
A direct substitution of $v=(9u+1-X)/16$ into \eqref{eq:Ruv} and \eqref{eq:Puv} gives
\begin{equation}\label{eq:S-Delta}
  S(u,X)=\frac{(u-1)^2\Delta}{1024}.
\end{equation}
Define
\begin{equation}
  \mathcal C:=\R[u,X,J]/(J^2+S(u,X)),
  \qquad
  \overline u=u,
  \quad
  \overline X=X,
  \quad
  \overline J=-J.
\end{equation}

\begin{theorem}[CUE $N=3$ trace splitting]\label{thm:N3-trace}
The algebra $\mathcal C$ is a free rank-two $\R[u,X]$-module
\begin{equation}
  \mathcal C=\R[u,X]\oplus \R[u,X]J.
\end{equation}
Define
\begin{equation}
  \langle A,B\rangle_{\mathcal C}:=\frac12\Tr_{\mathcal C/\R[u,X]}(A\overline B).
\end{equation}
For
\begin{equation}
  A=a+bJ,
  \qquad
  B=c+dJ,
\end{equation}
one has
\begin{equation}
  \langle A,B\rangle_{\mathcal C}=ac+bdS.
\end{equation}
In particular $\langle1,J\rangle_{\mathcal C}=0$, and
\begin{equation}
  \langle a+bJ,a+bJ\rangle_{\mathcal C}=a^2+Sb^2.
\end{equation}
On the chamber $\Delta>0$ and $u\ne1$, this form is positive definite.
\end{theorem}

\begin{proof}
Since $J^2=-S$ is a monic quadratic relation, $\mathcal C$ is free with basis $1,J$.  The involution is well-defined because $(-J)^2+S=J^2+S$.  The conjugates of $J$ are $J$ and $-J$, so
\begin{equation}
  \Tr(1)=2,
  \qquad
  \Tr(J)=0.
\end{equation}
Now
\begin{equation}
  (a+bJ)\overline{(c+dJ)}=(a+bJ)(c-dJ)=ac-bdJ^2+(bc-ad)J.
\end{equation}
Using $J^2=-S$, this is
\begin{equation}
  ac+bdS+(bc-ad)J.
\end{equation}
Taking half the trace gives $ac+bdS$.  Finally, \eqref{eq:S-Delta} gives $S>0$ on $\Delta>0$ and $u\ne1$.
\end{proof}

\begin{remark}
The $N=3$ compact Mellin computation proves that integer kernels are Hermitian-square sums on the phase cover and descend to quotient polynomials in $\Q[u,v]$ through this quadratic trace recurrence.  Already the descended quotient kernel $\Psi_2^\sharp$ is not a real sum of squares on the base, although its pullback to the phase cover is a Hermitian-square sum.  The stable theorem-level statement is therefore a quotient/lift obstruction: positivity lives naturally on the cover and need not descend as an SOS certificate on the quotient.
\end{remark}

\section{Weighted two-sector defects and $A$-lattices}\label{sec:defect-module}

For $m\in\R$, define
\begin{equation}
  \mathfrak D_2(m)=\R e_+\oplus\R e_-,
  \qquad
  \epsilon_m(ae_++be_-)=-m(a+b).
\end{equation}
Let
\begin{equation}
  e_{\vis}=e_++e_-,
  \qquad
  e_{\hid}=e_+-e_-.
\end{equation}
Then
\begin{equation}
  \epsilon_m(e_+)=-m,
  \qquad
  \epsilon_m(e_-)=-m,
\end{equation}
and
\begin{equation}
  \epsilon_m(e_{\vis})=-2m,
  \qquad
  \epsilon_m(e_{\hid})=0.
\end{equation}
For the peeled CUE measure,
\begin{equation}
  m=m_N=\frac{(N+1)(2N+1)}{6N}.
\end{equation}
Thus the primitive lift coefficient at $t=1$ is $-m_N$, the visible doubled finite defect is $-2m_N$, and the hidden difference sector is in the defect kernel.

More generally, define
\begin{equation}
  D_r=\Z^r,
  \qquad
  \epsilon_r(a_1,\dots,a_r)=-(a_1+\cdots+a_r).
\end{equation}
Then
\begin{equation}
  \ker(\epsilon_r)=\{(a_1,\dots,a_r):\sum_i a_i=0\}.
\end{equation}

\begin{theorem}[$A_{r-1}$ cancellation lattice]\label{thm:A-lattice}
For $r\ge2$,
\begin{equation}
  \ker(\epsilon_r)=\Z^r_0
\end{equation}
is the $A_{r-1}$ root lattice.  For $r=1$, the kernel is the zero lattice $A_0$.
\end{theorem}

\begin{proof}
The vectors
\begin{equation}
  \alpha_i=e_i-e_{i+1},\qquad 1\le i\le r-1,
\end{equation}
span $\Z^r_0$.  Their Gram matrix is
\begin{equation}
  \langle\alpha_i,\alpha_j\rangle=
  \begin{cases}
    2,&i=j,\\
    -1,&|i-j|=1,\\
    0,&|i-j|>1,
  \end{cases}
\end{equation}
which is the Cartan matrix of type $A_{r-1}$.  For $r=1$ the condition $a_1=0$ gives the zero lattice.
\end{proof}

This lattice is the diagonal augmentation skeleton.  It should not be confused with the full finite phase-space algebra, whose nonidentity sectors span $\mathfrak{sl}_r(\C)$ by \cref{thm:phase-space-basis}.

\section{Endpoint $t=1$, Stone defect, and coprime kernel}\label{sec:stone}

Let
\begin{equation}
  F_1(s)=\zeta(s)-1=\sum_{n\ge1}b(n)n^{-s},
  \qquad
  b(1)=0,
  \quad
  b(n)=1\ (n\ge2).
\end{equation}
For coprime $m,n$, define
\begin{equation}
  \Delta F_1(m,n):=b(mn)-b(m)b(n).
\end{equation}

\begin{theorem}[Boundary-supported multiplicativity defect]\label{thm:stone-boundary}
For $(m,n)=1$,
\begin{equation}
  \Delta F_1(m,n)=
  \begin{cases}
    1,&\text{exactly one of }m,n\text{ equals }1,\\
    0,&\text{otherwise}.
  \end{cases}
\end{equation}
Consequently
\begin{equation}
  \sum_{(m,n)=1}\Delta F_1(m,n)m^{-s}n^{-t}=\zeta(s)+\zeta(t)-2.
\end{equation}
\end{theorem}

\begin{proof}
If $m=n=1$, then $b(mn)=b(1)=0$ and $b(m)b(n)=0$.  If exactly one of $m,n$ equals $1$ and the other is at least $2$, then $b(mn)=1$ while $b(m)b(n)=0$.  If both are at least $2$, then $mn\ge2$ and $b(mn)=b(m)b(n)=1$.  Summing over coprime pairs gives two boundary sums:
\begin{equation}
  \sum_{n\ge2}n^{-t}+\sum_{m\ge2}m^{-s}=\zeta(t)-1+\zeta(s)-1.
\end{equation}
\end{proof}

Let
\begin{equation}
  u_0=\mathbf 1_{\{1\}},
  \qquad
  v_0=\mathbf 1_{\N_{\ge2}}.
\end{equation}
The Stone space of the clopen partition $\N=\{1\}\sqcup\N_{\ge2}$ has two points.  Its multiplication satisfies
\begin{equation}
  \xi_1\xi_1=\xi_1,
  \qquad
  \xi_1\xi_*=\xi_*\xi_1=\xi_*,
  \qquad
  \xi_*\xi_*=\xi_*.
\end{equation}
The pullback comultiplication gives
\begin{equation}
  \Delta(u_0)=u_0\otimes u_0,
  \qquad
  \Delta(v_0)=v_0\otimes u_0+u_0\otimes v_0+v_0\otimes v_0.
\end{equation}
Thus
\begin{equation}
  \Delta(v_0)-v_0\otimes v_0=v_0\otimes u_0+u_0\otimes v_0,
\end{equation}
which is exactly the unit-boundary defect from \cref{thm:stone-boundary}.

For the interior coprime part, define
\begin{equation}
  C(s,t):=\sum_{\substack{m,n\ge2\\(m,n)=1}}m^{-s}n^{-t}.
\end{equation}
Then
\begin{equation}\label{eq:Cst}
  C(s,t)=\frac{\zeta(s)\zeta(t)}{\zeta(s+t)}-\zeta(s)-\zeta(t)+1.
\end{equation}
If the nonunit coefficient is $\lambda v_0$, then for $m,n\ge2$ coprime the multiplicativity defect is
\begin{equation}
  \lambda-\lambda^2=\delta(\lambda).
\end{equation}
For the doubled sheet $\lambda=2$, one has
\begin{equation}
  \delta(2)=-2,
\end{equation}
and the interior defect is $-2C(s,t)$.

\section{Number-field shifted sheets}\label{sec:number-field}

Let $K$ be a number field.  Let $a_K(n)$ be the number of nonzero ideals of $\mathcal O_K$ of norm $n$, so
\begin{equation}
  \zeta_K(s)=\sum_{n\ge1}a_K(n)n^{-s}.
\end{equation}
For $|t|<1$, define
\begin{equation}
  \zeta_{K,t}(s):=\sum_{n\ge1}a_K(n)(n+t)^{-s}
  =\sum_{0\ne J\triangleleft\mathcal O_K}(N(J)+t)^{-s}.
\end{equation}

\begin{theorem}[Dedekind shifted deformation]\label{thm:Dedekind-shift}
For $|t|<1$,
\begin{equation}
  \zeta_{K,t}(s)=\sum_{m\ge0}(-1)^m\frac{(s)_m}{m!}t^m\zeta_K(s+m)
\end{equation}
for $\Re(s)>1$, and this identity gives a meromorphic continuation to the whole plane.  The function $\zeta_{K,t}$ has the same unique simple pole at $s=1$ as $\zeta_K$, with the same residue.
\end{theorem}

\begin{proof}
Apply \cref{thm:shifted-flow} to the Dirichlet series $L=\zeta_K$.  The standard meromorphic continuation and pole structure of $\zeta_K$ then carry over termwise in the displayed expansion.
\end{proof}

\begin{theorem}[Ideal-Dirichlet ratio kernel]\label{thm:Dedekind-ratio}
For every integer $r\ge1$, the ratio
\begin{equation}
  Q_{K,r}(s):=\frac{\zeta_K(s+r)}{\zeta_K(s)}
\end{equation}
has an ideal-Dirichlet expansion
\begin{equation}
  Q_{K,r}(s)=\sum_{\mathfrak a\ne0}\alpha_{K,r}(\mathfrak a)N(\mathfrak a)^{-s},
\end{equation}
where $\alpha_{K,r}$ is multiplicative and
\begin{equation}
  \alpha_{K,r}(\mathfrak p^k)=-\frac{N(\mathfrak p)^r-1}{N(\mathfrak p)^{rk}}
  \qquad(k\ge1).
\end{equation}
Equivalently,
\begin{equation}
  \alpha_{K,r}(\mathfrak a)=(-1)^{\omega(\mathfrak a)}
  \frac{\prod_{\mathfrak p\mid\mathfrak a}(N(\mathfrak p)^r-1)}{N(\mathfrak a)^r}.
\end{equation}
\end{theorem}

\begin{proof}
Using the Euler product,
\begin{equation}
  \zeta_K(s)=\prod_{\mathfrak p}(1-N(\mathfrak p)^{-s})^{-1}.
\end{equation}
Thus
\begin{equation}
  Q_{K,r}(s)=\prod_{\mathfrak p}
  \frac{1-N(\mathfrak p)^{-s}}{1-N(\mathfrak p)^{-s-r}}.
\end{equation}
For $q=N(\mathfrak p)$ and $x=q^{-s}$, the local factor is
\begin{equation}
  \frac{1-x}{1-q^{-r}x}=1-\sum_{k\ge1}(1-q^{-r})q^{-r(k-1)}x^k.
\end{equation}
Thus the coefficient at $\mathfrak p^k$ is
\begin{equation}
  -(1-q^{-r})q^{-r(k-1)}=-\frac{q^r-1}{q^{rk}},
\end{equation}
and multiplicativity follows from the Euler product.
\end{proof}

\begin{corollary}\label{cor:Dedekind-nonEuler}
If $0<t<1$, then $\zeta_{K,t}$ is not an ordinary Dirichlet series in the variable $n^{-s}$ and has no classical Euler product.
\end{corollary}

\begin{proof}
The first nonzero coefficient of $\zeta_K$ occurs at $n=1$, because exactly one ideal has norm $1$, namely $\mathcal O_K$.  The leading exponential scale of $\zeta_{K,t}(\sigma)$ as $\sigma\to+\infty$ is therefore $(1+t)^{-\sigma}$, which cannot equal $n_0^{-\sigma}$ for an integer $n_0$ when $0<t<1$.  This is the same argument as \cref{thm:non-eulerianity}.
\end{proof}

\section{Consolidated structural theorem}\label{sec:synthesis}

\begin{theorem}[Consolidated package]\label{thm:consolidated}
The constructions above fit into the following exact algebraic package.
\begin{enumerate}[label=\textup{(\arabic*)},leftmargin=2.4em]
\item The split-zero semiring separates arithmetic value from support:
\begin{equation}
  p_R^{-1}(0_R)=\{\tau,0_R\},
  \qquad
  \chi_R^{-1}(0)=\{\tau\},
  \qquad
  \chi_R(0_R)=1.
\end{equation}
\item The linear theory over $G(R)$ is not ordinary module theory; it is the diagram category
\begin{equation}
  G(R)\text{-}\SMod\simeq\Diag_R,
\end{equation}
with cancellative sector $R$-Mod and free objects given by Boolean cubes of $R$-modules.
\item The quotient/lift principal parts
\begin{equation}
  q=1/u,
  \qquad
  \ellp=v/u
\end{equation}
have function parities $+$ and $-$, respectively, but residues reverse parity:
\begin{equation}
  q\mapsto -e_0+e_1,
  \qquad
  \ellp\mapsto e_0+e_1.
\end{equation}
\item The completed shifted-Hurwitz principal part is
\begin{equation}
  \PP_{\{0,1\}}\Lambda_{\mu,t}
  =\bigl(1+t m_1(\mu)\bigr)q-tm_1(\mu)\ellp.
\end{equation}
The finite two-end defect is carried by the anti-invariant lift coefficient and is canceled by projective infinity.
\item For the peeled CUE measure,
\begin{equation}
  m_1(\mu_N)=\frac{(N+1)(2N+1)}{6N}.
\end{equation}
Hence the primitive CUE lift coefficient is
\begin{equation}
  -t\frac{(N+1)(2N+1)}{6N},
\end{equation}
and for $N=3$ it is $-14t/9$.
\item Fixed-$N$ CUE derivative moments form a compact Mellin stack: the positive-integer lattice reconstructs the compact measure by Bernstein-Hausdorff inversion, and the finite Bernstein-Hurwitz approximants preserve zero counts in stable regions.
\item For every finite abelian group $A$ of order $r$, the $r^2-1$ nonidentity sectors of $A\times\widehat A$ form an orthogonal operator basis of $\mathfrak{sl}_r(\C)$.  The root lattice $A_{r-1}$ is the diagonal augmentation skeleton, not the full phase-space double.
\item The first-derivative $s$-fold CUE microscopic exponent satisfies
\begin{equation}
  s^2+2s=(s+1)^2-1=\dim\mathfrak{sl}_{s+1}(\C),
\end{equation}
and is realized by the finite phase-space basis for $A=\Z/(s+1)\Z$.
\item The first nontrivial $N=3$ CUE quotient/lift obstruction is modeled by a quadratic trace algebra
\begin{equation}
  \mathcal C=\R[u,X]\oplus\R[u,X]J,
  \qquad
  J^2=-S,
  \qquad
  \overline J=-J,
\end{equation}
with trace form $\langle a+bJ,a+bJ\rangle=a^2+Sb^2$.
\end{enumerate}
\end{theorem}

\begin{proof}
Item (1) is \cref{prop:DR,thm:unique-boolean-character}.  Item (2) is \cref{thm:semimodule-classification,cor:cancellative-sector,thm:free-boolean-cube}.  Item (3) is \cref{thm:residue-reversal}.  Item (4) is \cref{thm:pp-decomp,cor:lift-defect-half}.  Item (5) is \cref{thm:first-CUE-moment,cor:CUE-pp}.  Item (6) is \cref{thm:compact-support,thm:bernstein,thm:finite-zero-stack}.  Item (7) is \cref{thm:phase-space-basis,cor:A-root-lattice}.  Item (8) is \cref{thm:first-jet-block,cor:CUE-phase-exponent}.  Item (9) is \cref{thm:N3-trace}.
\end{proof}

\subsection*{Final dictionary}
The structure is therefore not a naive product of unrelated $\Z/2$ effects.  It is a layered package:
\begin{equation}
  \boxed{\text{idempotent Boolean support thickening }G(R)}
\end{equation}
followed by
\begin{equation}
  \boxed{\text{ramified quadratic quotient/lift cover }\C(u)\subset\C(u)(v)}
\end{equation}
and finally by
\begin{equation}
  \boxed{\text{finite CUE phase-space enlargement }A\times\widehat A.}
\end{equation}
The support layer explains why $\tau$ is not the supported zero.  The quadratic layer explains why residues reverse function parity.  The CUE-Hurwitz layer explains why the finite endpoint defect is governed by the first compact Mellin moment.  The phase-space layer explains why the first-derivative exponent $s^2+2s$ belongs naturally to the full traceless algebra $\mathfrak{sl}_{s+1}$ rather than only to an $A_s$ diagonal lattice.

\begin{thebibliography}{99}

\bibitem{GMS}
A. Grover, F. Mezzadri, and N. Simm,
\emph{Higher order derivative moments of CUE characteristic polynomials and the Riemann zeta function}, arXiv:2604.03051, 2026.

\bibitem{CompactMellin}
\emph{Compact Mellin peeling, transform algebras, and the first noninteger layers of CUE derivative moments}, project note, April 2026.

\bibitem{SplitZeroSecondary}
\emph{A secondary note on split-zero globalization: semimodule classification, multiplicative reconstruction, and all-orders zeta deformation}, project note, April 2026.

\bibitem{KeatingSnaith}
J. P. Keating and N. C. Snaith,
\emph{Random matrix theory and $\zeta(1/2+it)$}, Comm. Math. Phys. 214 (2000), 57--89.

\bibitem{ConreyRubinsteinSnaith}
J. B. Conrey, M. O. Rubinstein, and N. C. Snaith,
\emph{Moments of the derivative of characteristic polynomials with an application to the Riemann zeta function}, Comm. Math. Phys. 267 (2006), 611--629.

\bibitem{Dehaye}
P.-O. Dehaye,
\emph{Joint moments of derivatives of characteristic polynomials}, Algebra Number Theory 2 (2008), 31--68.

\bibitem{Hughes}
C. P. Hughes,
\emph{On the characteristic polynomial of a random unitary matrix and the Riemann zeta function}, Ph.D. thesis, University of Bristol, 2001.

\bibitem{Macdonald}
I. G. Macdonald,
\emph{Symmetric Functions and Hall Polynomials}, Oxford University Press, 1995.

\bibitem{DLMF}
NIST Digital Library of Mathematical Functions, Chapter 25: Zeta and related functions, especially Section 25.11 on the Hurwitz zeta function.

\bibitem{MilneANT}
J. S. Milne,
\emph{Algebraic Number Theory}, available from the author's course notes.

\bibitem{Neukirch}
J. Neukirch,
\emph{Algebraic Number Theory}, Grundlehren der mathematischen Wissenschaften 322, Springer, 1999.

\bibitem{DrosteKuichVogler}
M. Droste, W. Kuich, and H. Vogler, eds.,
\emph{Handbook of Weighted Automata}, Springer, 2009.

\bibitem{Thas}
K. Thas, ed.,
\emph{Absolute Arithmetic and $F_1$-Geometry}, European Mathematical Society, 2016.

\end{thebibliography}

\end{document}
